\section{Introducción}

\subsection*{Propósito}

Este documento evalúa el marco ético del proyecto \textbf{Astro-Denoise}, una iniciativa de investigación científica que busca crear herramientas de acceso abierto para mejorar la calidad de imágenes astronómicas. El análisis se enfoca en cuatro principios fundamentales: reproducibilidad científica, equidad en acceso a datos, transparencia operativa y responsabilidad institucional.

\subsection*{Hallazgos Clave}

El proyecto Astro-Denoise está diseñado como respuesta ética a una brecha sistémica en el acceso a pipelines avanzados de procesamiento de datos astronómicos. Instituciones en países fuera de la alianza oficial LSST (Large Synoptic Survey Telescope) enfrentan limitaciones para participar activamente en la ciencia astronómica moderna. Este proyecto mitiga esa brecha mediante:

\begin{itemize}
    \item Código abierto y reproducible accesible globalmente
    \item Uso de datos públicos simulados (LSST DC2) disponibles sin restricción geográfica
    \item Documentación clara en español, facilitando inclusión de investigadores de habla hispana
    \item Herramientas modernas que no requieren infraestructura costosa de cómputo especializado
\end{itemize}

\subsection*{Evaluación del Marco Ético}

Se ha identificado la aplicación exitosa de tres principios éticos clave:

\begin{enumerate}
    \item \textbf{Reproducibilidad y Transparencia Científica} (\cite{unesco2021}): El proyecto garantiza que todo resultado sea auditável y reproducible mediante datos versionados, código abierto y reportes estructurados.
    
    \item \textbf{Equidad en Acceso a Recursos Científicos} (\cite{minciencias2022}): Astro-Denoise reduce la brecha geográfica e institucional que limita la participación científica, alineándose con políticas nacionales de ciencia abierta.
    
    \item \textbf{Gobernanza Responsable de Datos} (\cite{wilkinson2016}): El manejo de datos sigue principios FAIR (Accesible, Interoperable, Reutilizable) y respeta la Ley 1581 de Protección de Datos.
\end{enumerate}

\subsection*{Riesgos Identificados y Mitigaciones}

Se han identificado dos riesgos éticos principales con estrategias de mitigación implementadas:

\begin{table}[H]
\centering
\small
\begin{tabular}{|p{3cm}|p{3cm}|p{3cm}|}
\hline
\textbf{Riesgo Ético} & \textbf{Descripción} & \textbf{Mitigación} \\
\hline
Reprodu Limitada & Resultados científicos no auditables & Versionamiento de datos, código abierto, reportes estructurados \\
\hline
Acceso Desigual & Instituciones sin recursos no pueden validar & Herramientas sin requerimientos extremos, documentación pública \\
\hline
\end{tabular}
\end{table}

\subsection*{Recomendación}

El proyecto \textbf{se alinea con los criterios éticos esperados} para investigación científica responsable en el contexto colombiano e internacional. Se recomienda mantener el compromiso de acceso abierto y reproducibilidad a lo largo de su ciclo de vida.
